\abstracttext{As audiências públicas da Câmara dos Deputados chegam ao público sobretudo por matérias
jornalísticas que selecionam poucas opiniões, sem indicar onde cada uma aparece na transcrição. Este
trabalho investiga se um perfil escrito só com as falas anteriores de uma pessoa ajuda um modelo de
linguagem a reconhecer o que a imprensa lhe atribuiu numa audiência posterior. A Unidade Deliberativa
Verificável localiza na fala da própria pessoa evidência candidata para 2.105 das 2.203 opiniões do
PublicHearingBR, com a localização na transcrição, o tipo e a origem do vínculo. Os perfis por ator,
escritos por um modelo exclusivamente com as falas de cada pessoa nas audiências de treino, têm
blocos de posições, critérios e valores, alinhamentos declarados e forma de argumentar. Na simulação,
um modelo responde como a pessoa a partir do perfil e de trechos recuperados das falas anteriores. A
avaliação usa como gabarito as
opiniões publicadas sobre audiências posteriores, em perguntas de múltipla escolha com distratores da
mesma audiência.
<Resultados principais: acerto de cada condição no teste e diferenças entre condições com intervalo
de confiança.>}
